\section{Interpretazione Astratta in LiSA}

\subsection{Cos'è LiSA}

\begin{definizione}{LiSA (Library for Static Analysis)}{lisa_def}
LiSA è una libreria open-source scritta in Java, sviluppata e mantenuta dal
gruppo di ricerca SSV dell'Università Ca' Foscari di Venezia.
Il suo scopo è fornire un framework modulare e estensibile per la costruzione
di analizzatori statici.
\end{definizione}

I principi fondamentali su cui si basa LiSA sono:
\begin{itemize}
    \item Un programma in LiSA è rappresentato come un insieme di control
    flow graphs (CFGs), più eventuali informazioni ausiliarie.
    \item I CFG codificano il flusso di controllo nella loro struttura,
    astraendo dai costrutti sintattici del linguaggio sorgente.
    \item I nodi del CFG sono specifici del linguaggio, ma vengono riscritti
    in \textbf{espressioni simboliche} generiche.
\end{itemize}

\begin{nota}{Espressioni simboliche}{symbolic_expressions}
    La riscrittura dei CFG in espressioni simboliche generiche permette di
    separare la sintassi dalla semantica, consentendo ai domini astratti di
    operare in modo agnostico rispetto al linguaggio.
\end{nota}
\subsection{Architettura di LiSA}

L'architettura di LiSA è progettata per essere altamente modulare, separando
i componenti dell'analisi per massimizzare la riusabilità.

\begin{figure}[H]
    \centering
    \includegraphics[width=0.9\textwidth]{images/th_11/01.png}
    \caption{L'immagine illustra la visione d'insieme dell'architettura di
    LiSA.}
    \label{fig:th_11_01}
\end{figure}

I componenti principali, come mostrato in \Cref{fig:th_11_01}, sono:
\begin{itemize}
    \item \textbf{Li front-end}: Traduce il programma sorgente ($P_i$) in una
    collezione di CFG.
    \item \textbf{LiSA ``core''}: È il cuore dell'analizzatore e include:
    \begin{itemize}
        \item \textbf{CFG fixpoint}: L'engine che implementa l'algoritmo di
        punto fisso basato su worklist per calcolare la semantica astratta.
        \item \textbf{Statement semantics}: Logica per tradurre i nodi del
        CFG in espressioni simboliche.
        \item \textbf{Domain}: Il dominio astratto che definisce le
        proprietà da tracciare (es. intervalli, segni).
        \item \textbf{Memory handling}: Componente che modella la
        memoria (heap).
        \item \textbf{Program fixpoint}: Gestisce l'analisi interprocedurale
        e la risoluzione delle chiamate (call resolution).
    \end{itemize}
    \item \textbf{Checks}: Componenti che verificano la presenza di specifici
    bug o proprietà sui risultati dell'analisi.
    \item \textbf{Analysis dump}: Produce report e warning per l'utente.
\end{itemize}

\subsection{Calcolo del punto fisso sul CFG}

L'analisi intra-procedurale in LiSA si basa su un calcolo di punto
fisso eseguito direttamente sul CFG.

\begin{definizione}{Algoritmo di punto fisso del CFG}{cfg_fixpoint_algo}
L'algoritmo è basato su un approccio worklist per raggiungere la stabilità.
Lo pseudocodice è il seguente:
\begin{verbatim}
result = {}
ws = {entry_node}
while (ws is not empty) do
    st = ws.pop()
    in_st = LUB { out_m | m in preds(st) }
    out_st = semantics(st, in_st)
    if (not out_st <= result[st]) then
        if (widening_threshold_reached) then
            result[st] = result[st] WIDEN out_st
        else
            result[st] = result[st] LUB out_st
        end if
        ws.push(successors(st))
    end if
end while
\end{verbatim}
\end{definizione}

\begin{nota}{Componenti dell'algoritmo}{fixpoint_components}
L'algoritmo si basa su operazioni definite dal dominio astratto, che deve
essere un reticolo:
\begin{itemize}
    \item \textbf{LUB (least upper bound)}: Unisce l'informazione
    proveniente da più cammini di controllo.
    \item \textbf{Partial order}: Verifica se un nuovo stato aggiunge
    informazione rispetto a quello già calcolato, per determinare se la
    computazione si è stabilizzata.
    \item \textbf{Widening}: Operatore di allargamento (prossimi capitoli) che
    garantisce la terminazione dell'analisi su reticoli di altezza infinita,
    sacrificando la precisione per forzare la convergenza.
\end{itemize}
Gli stati calcolati devono quindi essere elementi di un reticolo
(\Cref{def:lattice}).
\begin{figure}[H]
    \centering
    \includegraphics[width=0.3\textwidth]{images/th_11/02.png}
    \caption{L'immagine illustra le operazioni principali su un reticolo.}
    \label{fig:th_11_02}
\end{figure}
\end{nota}

\subsection{Semantica degli statement}

In LiSA, la semantica di uno statement astrae dalle operazioni concrete e si concentra
su tre funzionalità principali che un dominio deve supportare.

\begin{definizione}{SemanticDomain}{semantic_domain_def}
Uno stato astratto in LiSA deve implementare
l'interfaccia \texttt{SemanticDomain}, che fornisce metodi per:
\begin{itemize}
    \item \textbf{assign(id, expr, \ldots)}: Gestisce l'assegnamento di
    un'espressione a una variabile.
    \item \textbf{smallStepSemantics(expr, \ldots)}: Valuta un'espressione
    simbolica.
    \item \textbf{satisfies(expr, \ldots)}: Verifica se un'espressione è
    soddisfatta dallo stato astratto.
    \item \textbf{assume(expr, \ldots)}: Raffina lo stato astratto assumendo
    che un'espressione sia vera.
\end{itemize}

\begin{figure}[H]
    \centering
    \includegraphics[width=0.6\textwidth]{images/th_11/03.png}
    \caption{L'immagine illustra le operazioni principali sul
    \texttt{SemanticDomain}.}
    \label{fig:th_11_03}
\end{figure}

\end{definizione}

\begin{esempio}{Semantica di un condizionale}{cond_semantics_ex}
La semantica di un'istruzione condizionale \texttt{if (cond) then \ldots}
viene gestita come segue:
\begin{verbatim}
// Calcola lo stato per il ramo "then"
tbranch = domain.assume(cond).smallStepSemantics(left)
// Calcola lo stato per il ramo "else"
fbranch = domain.assume(cond.negate()).smallStepSemantics(right)

// Verifica la soddisfacibilità per potare i rami
if domain.satisfies(cond) == SATISFIED:
    return tbranch
else if domain.satisfies(cond) == NOT_SATISFIED:
    return fbranch
else: // Se non so, unisco i risultati
    return tbranch.lub(fbranch)
\end{verbatim}

\end{esempio}

\subsection{Gestione della memoria e stato astratto}

Per garantire la modularità, LiSA separa la gestione dei valori dalla
gestione della memoria.

\subsubsection{Gestione della memoria dinamica}

\begin{nota}{Indipendenza dalla memoria}{memory_independence}
I domini astratti (es. intervalli, segni) non dovrebbero preoccuparsi di
come la memoria è organizzata.
L'idea è trattare ogni locazione di memoria (es. un campo di un oggetto)
come una ``variabile fittizia'' (fake variable).
\end{nota}

\begin{esempio}{Accesso a un campo di un oggetto}{field_access_ex}
    Considerando il seguente frammento di codice che alloca due oggetti:
\begin{verbatim}
class A {
    int f, g;
}
x = new A();
y = new A();
\end{verbatim}
L'Heap Domain mappa ogni oggetto e i suoi campi a una locazione di memoria
concreta, che viene trattata come una variabile fittizia.
Come mostrato in \ref{fig:th_11_04}, l'oggetto puntato da \texttt{x} si
trova alla locazione $L_1$, mentre quello puntato da \texttt{y} si trova
in $L_2$.
Di conseguenza, l'accesso al campo \texttt{f} di \texttt{x} (\texttt{x.f})
viene risolto nella variabile fittizia $L_{1.1}$, mentre \texttt{y.f} viene
risolto in $L_{2.1}$.
Sarà poi compito del Value Domain associare un valore astratto a queste
variabili fittizie (es. $L_{1.1} \to [5, 5]$).

\begin{figure}[H]
    \centering
    \includegraphics[width=0.6\textwidth]{images/th_11/04.png}
    \caption{L'immagine illustra la modellazione della memoria dinamica. Ogni
    oggetto e campo viene mappato a una locazione di memoria unica.}
    \label{fig:th_11_04}
\end{figure}

\end{esempio}

\subsubsection{Lo stato astratto}

Lo stato astratto (\ref{fig:th_11_05}) in LiSA è tipicamente composto da due parti:
\begin{itemize}
    \item \textbf{Heap domain}: Si occupa di modellare la memoria.
    La sua responsabilità è risolvere un'espressione che accede alla memoria
    (es. \texttt{x.f}) nella corrispondente variabile fittizia che
    rappresenta quella locazione.
    \item \textbf{Value domain}: Mappa le variabili (sia quelle dello stack
    che quelle fittizie dello heap) ai loro valori astratti (es. intervalli,
    segni).
\end{itemize}

\begin{figure}[H]
    \centering
    \includegraphics[width=0.4\textwidth]{images/th_11/05.png}
    \caption{L'immagine illustra la composizione dello stato astratto in LiSA.}
    \label{fig:th_11_05}
\end{figure}

A seconda della precisione richiesta, l'heap domain può modellare le locazioni
di memoria in modi diversi (es. context-sensitive, context-insensitive),
influenzando il trade-off tra precisione e performance.

\subsection{Analisi interprocedurale}

L'analisi interprocedurale gestisce le chiamate a funzione, orchestrando
le analisi intra-procedurali dei singoli CFG.

\begin{definizione}{Analisi Interprocedurale}{interproc_analysis}
Questo componente è responsabile di:
\begin{itemize}
    \item Analizzare i CFG in isolamento.
    \item Iniziare l'analisi dal ``main'' e seguire le chiamate a funzione.
    \item Utilizzare un \textbf{call graph} per risolvere dinamicamente i
    target di una chiamata (es. in presenza di polimorfismo).
\end{itemize}
\end{definizione}

\begin{figure}[H]
    \centering
    \includegraphics[width=0.4\textwidth]{images/th_11/06.png}
    \caption{L'immagine illustra l'analisi interprocedurale in LiSA.}
    \label{fig:th_11_06}
\end{figure}


\subsection{Domini astratti: relazionali e non-relazionali}

I domini di valori si dividono in due categorie principali:
\begin{itemize}
    \item \textbf{Domini non-relazionali}: Tracciano i valori di ogni variabile
    in modo indipendente.
    \item \textbf{Domini relazionali}: Tracciano le relazioni tra i valori di
    più variabili.
\end{itemize}

\subsubsection{Domini non-relazionali}

\begin{definizione}{Dominio non-relazionale}{non_rel_domain}
Un dominio non-relazionale traccia i valori di ogni variabile in modo
indipendente, tramite una mappa da identificatori a valori astratti.

In questo tipi di domini la logica di valutazione è specifica per il dominio.
\end{definizione}

\begin{esempio}{Esempi di domini non-relazionali}{non_rel_domain_ex}
Alcuni esempi di domini non-relazionali sono:
\begin{itemize}
    \item Constant Propagation ($x \to 3$).
    \item Intervals ($x \to [2, 10]$).
\end{itemize}
\end{esempio}

\begin{figure}[H]
    \centering
    \includegraphics[width=0.6\textwidth]{images/th_11/07.png}
    \caption{L'immagine illustra le classi di domini non-relazionali in LiSA.}
    \label{fig:th_11_07}
\end{figure}


\subsubsection{Domini relazionali}

\begin{definizione}{Dominio relazionale}{rel_domain}
Un dominio relazionale traccia le relazioni tra i valori di più variabili.
\begin{itemize}
    \item \textbf{Esempi}: Upper bounds ($x > y$), Linear inequalities ($x+y > 10, z < 0$).
\end{itemize}
\end{definizione}


\begin{esempio}{Esempi di domini relazionali}{rel_domain_ex}
Alcuni esempi di domini relazionali sono:
\begin{itemize}
    \item Upper bounds ($x > y$).
    \item Linear inequalities ($x+y > 10, z < 0$).
\end{itemize}
\end{esempio}